\bookchapter{Capítulo 59: Conhecendo o iOS — o ecossistema da Apple e o ecossistema de desenvolvimento}

\chapterintro{Parabéns por tomar essa decisão — passar de “quem usa o iPhone” para “quem cria apps para o iPhone”. Antes de colocar a mão na massa, vamos estender o mapa na mesa: de onde vem o iOS? O que o “jardim murado” da Apple significa para quem desenvolve? O que acontece, do momento em que você digita o código até o app chegar ao celular do usuário?}

\booksection{A origem do iOS: um celular que mudou uma indústria}

Em janeiro de 2007, Steve Jobs tirou o primeiro iPhone do bolso em uma apresentação. Naquela época, ele nem aceitava apps de terceiros — só em 2008 a Apple lançou a App Store e o iPhone SDK, permitindo que desenvolvedores do mundo todo criassem software para o iPhone. Foi nesse instante que nasceu a “economia de aplicativos”: hoje, a App Store tem milhões de apps e já repassou centenas de bilhões de dólares aos desenvolvedores.

O sistema também evoluiu: no começo se chamava iPhone OS e em 2010 passou a se chamar iOS; a numeração avançou até o iOS 17 (2023) e, em 2025, a Apple mudou os nomes para acompanhar o ano — o iOS 26 foi lançado em 2025 e, em agosto de 2026, a versão atual é o iOS 27. O nome pode mudar, mas uma coisa nunca mudou: a disposição dos usuários do iOS para pagar e a reputação de qualidade dos seus apps são reconhecidamente as melhores do ecossistema móvel — o mesmo app costuma faturar bem mais no iOS do que no Android.

\begin{center}
\begin{tabularx}{\linewidth}{|X|X|X|}
\hline
\textbf{Marco} & \textbf{Ano} & \textbf{O que foi importante para os desenvolvedores} \tabularnewline
\hline
Estreia da App Store & 2008 & Início do ecossistema de desenvolvimento de terceiros \tabularnewline
Lançamento do Swift & 2014 & Chegada da linguagem moderna para desenvolvimento iOS \tabularnewline
Lançamento do SwiftUI & 2019 & Framework de interface declarativa; o fio condutor deste livro \tabularnewline
iOS 26 (primeiro ano da nova nomenclatura) & 2025 & Nomes seguem o ano; integração completa do Apple Intelligence \tabularnewline
iOS 27 & 2026 & Versão mais recente (referência deste livro) \tabularnewline
\hline
\end{tabularx}
\end{center}

\begin{center}\small\color{gray}Linha do tempo do desenvolvimento iOS (até agosto de 2026)\end{center}

\booksection{O jardim murado: as “regras” e os “benefícios” do ecossistema da Apple}

O ecossistema da Apple costuma ser descrito como um “jardim murado” (Walled Garden) — muitas regras e controle rígido: os apps só entram na loja depois da revisão, o desenvolvimento exige um Mac e o Xcode, e a distribuição acontece principalmente pela App Store. Parece apertado? Mas, para quem desenvolve, essas regras rendem benefícios que valem dinheiro de verdade:

\begin{itemize}
  \item Fragmentação de aparelhos praticamente nula: os modelos de iPhone são pouquíssimos, com telas e versões de sistema bem padronizadas — adaptando seu app para 5 modelos, você alcança 95\% dos usuários (no Android, pode ser preciso adaptar para dezenas);
  \item Taxa de atualização do sistema altíssima: em poucas semanas depois do lançamento de um sistema novo, sete ou oito em cada dez usuários já atualizaram; dá para usar as novas APIs sem medo, sem precisar “cuidar de sistemas de cinco anos atrás”;
  \item Usuários acostumados a pagar: quem usa a App Store está mais disposto a gastar com apps e compras dentro do aplicativo;
  \item A revisão é rigorosa, mas a loja fica limpa: poucos apps falsos e maliciosos, e os bons produtos têm mais chance de serem vistos.
\end{itemize}

\begin{tcbwarn}{Atenção: falando primeiro sobre os custos}
Desenvolver para iOS tem exigências inegociáveis: é preciso ter um Mac (o Xcode só existe no macOS) e publicar na loja exige uma conta de desenvolvedor de 99 dólares por ano (a conta gratuita de desenvolvedor individual permite testar no aparelho localmente, mas não publicar). Pense bem nesse investimento antes de partir — a boa notícia é que, na fase de aprendizado, a conta gratuita dá e sobra; falaremos disso no Capítulo 60.

\end{tcbwarn}

\booksection{Em qual linguagem programar: Objective-C ou Swift?}

A história do desenvolvimento iOS também teve dois “protagonistas”. O Objective-C é o veterano fundador — usado de 2008 a 2014; o Swift é a linguagem nova lançada pela Apple em 2014 e que evoluiu até o Swift 6 (lançado em 2024, com checagem rigorosa de concorrência). Em 2026, praticamente todos os projetos novos usam Swift.

Uma dica sincera para quem parte do zero: aprenda direto Swift e nem precisa olhar o Objective-C — a não ser que, depois de contratado, você precise dar manutenção a código legado de dez anos atrás. As vantagens do Swift: sintaxe limpa (sem o fardo de ponto e vírgula e arquivos de cabeçalho da era C), o mecanismo de opcionais (Optional), que elimina no nível da linguagem aquele pesadelo da indústria chamado “crash de ponteiro nulo”, integração perfeita com o SwiftUI e o Swift Playgrounds no iPad para aprender brincando.

Estilo Objective-C (relíquia antiga)

\begin{lstlisting}
NSString *name = @"João"; [name stringByAppendingString:@"Olá"]; // colchetes + ponteiro + ponto e vírgula: o trio clássico
\end{lstlisting}

Estilo Swift (moderno e limpo)

\begin{lstlisting}
var name = "João" name += "Olá" // direto como uma frase em linguagem natural
\end{lstlisting}

\booksection{Como um app nasce?}

Como em toda plataforma, um app iOS também nasce numa linha de produção. Só que a versão da Apple tem duas “inspeções de segurança” a mais:

Ideia

Design (interface + funções)

Escrever código (Swift + SwiftUI)

Compilar e assinar

Teste (emulador + aparelho real)

Revisão da Apple

Publicado na App Store

\begin{center}\small\color{gray}Figura 1: o fluxo completo de um app iOS, da ideia à loja (repare na etapa extra de “revisão”)\end{center}

Conceitos-chave do fluxo:

\begin{itemize}
  \item Assinatura (Signing): todo app iOS precisa ser “assinado” para poder ser instalado — é a prova de que “esse pacote veio de você e não foi alterado”. Na fase de aprendizado, o Xcode assina automaticamente com o seu Apple ID; para publicar, usa-se o certificado de desenvolvedor. Detalhes no Capítulo 67.
  \item Simulador (Simulator): um iPhone virtual rodando no Mac, o campo de batalha principal do desenvolvimento e dos testes. Como ele usa o processador do Mac, é extremamente fluido — uma experiência bem melhor que a do emulador do Android.
  \item TestFlight: a ferramenta oficial da Apple para distribuir versões de teste — antes de publicar, você envia a versão beta para até 10.000 testadores, que instalam e experimentam pelo próprio app TestFlight. É um capricho exclusivo do ecossistema iOS.
\end{itemize}

\booksection{Panorama da stack tecnológica do iOS: seu “arsenal”}

\begin{center}
\begin{tabularx}{\linewidth}{|X|X|X|}
\hline
\textbf{Camada} & \textbf{Escolha predominante (2026)} & \textbf{Para que serve} \tabularnewline
\hline
Linguagem & Swift 6 & Linguagem para escrever a lógica \tabularnewline
Interface & SwiftUI + diretrizes de design da Apple (HIG) & Desenhar a interface com código (o fio condutor deste livro) \tabularnewline
Arquitetura & MVVM + @Observable & Separar interface e lógica \tabularnewline
Armazenamento de dados & SwiftData (banco de dados) + UserDefaults (configurações) & Guardar dados no celular \tabularnewline
Rede & URLSession + async/await + Codable & Conversar com servidores \tabularnewline
Ferramenta de build & Xcode Build System / xcodebuild & Transformar código em pacote instalável \tabularnewline
Ambiente de desenvolvimento & Xcode (exclusivo do macOS) & IDE oficial da Apple \tabularnewline
\hline
\end{tabularx}
\end{center}

\begin{center}\small\color{gray}Visão rápida da stack iOS predominante em 2026\end{center}

Percebeu? Essa lista é bem mais curta que a do Android — porque a Apple “resolve” muita coisa por você. O preço de um ecossistema fechado é a liberdade; a recompensa é a tranquilidade: sem precisar decidir qual biblioteca de rede ou qual banco de dados usar, o que a Apple fornece costuma ser a melhor solução.

\begin{tcbpractice}{Mãos à obra: conheça o “cadastro do seu celular na Apple”}
Abra [Ajustes] > [Geral] > [Sobre] no iPhone e anote três coisas: a versão do iOS, o nome do modelo e o espaço disponível. Depois dê uma volta pela aba “Hoje” da App Store e observe as capturas de tela e os textos dos apps em destaque — a partir de hoje, seu olhar sobre a loja passa de “usuário” para “desenvolvedor”.

\end{tcbpractice}

\begin{tcbsummary}{Resumo do capítulo}
\begin{itemize}
  \item Em 2008 estreou a App Store e nasceu a “economia de apps”; em 2026, a versão atual é o iOS 27
  \item O balanço do jardim murado: muitas regras, mas poucos aparelhos, sistemas atualizados, usuários que pagam e uma loja limpa
  \item Exigências: é obrigatório ter um Mac; publicar exige conta de desenvolvedor de 99 dólares/ano; na fase de aprendizado, a conta gratuita basta
  \item Linguagem: vá direto para o Swift; o Objective-C fica só para código legado
  \item A linha de produção tem uma etapa a mais que a do Android: a revisão da Apple; o TestFlight é a ferramenta ideal para distribuir betas
  \item Stack padrão de 2026: Swift 6 + SwiftUI + MVVM + SwiftData + URLSession + Xcode
\end{itemize}

\end{tcbsummary}
